\documentclass[11pt,a4paper]{letter}

\usepackage[T1]{fontenc}
\usepackage[utf8]{inputenc}
\usepackage{lmodern}
\usepackage[margin=25mm]{geometry}
\usepackage{microtype}
\usepackage[hidelinks]{hyperref}

\hypersetup{
  pdftitle={Cover Letter: Reusable Neural Guidance for Petri Net Alignments with Replay Verification},
  pdfauthor={Alessandro Berti and Wil M. P. van der Aalst}
}
\pagestyle{empty}
\setlength{\longindentation}{0pt}
\setlength{\indentedwidth}{\textwidth}
\setlength{\parskip}{0.5em}
% The letter class leaves six medium skips above the typed signature.
\setlength{\medskipamount}{3pt}

\address{Alessandro Berti\\
  Process and Data Science (PADS)\\
  RWTH Aachen University\\
  Aachen, Germany\\
  \href{mailto:a.berti@pads.rwth-aachen.de}{a.berti@pads.rwth-aachen.de}}
\signature{Alessandro Berti\\
  Corresponding author, on behalf of both authors}
\date{\today}

% Compile from this directory with:
% pdflatex -output-directory=build -interaction=nonstopmode -halt-on-error cover_letter.tex
% Before submission, the corresponding author should verify the publication
% status and author-approval declarations below. If the manuscript was previously
% submitted to an MDPI journal, add that history and the previous manuscript ID.

\begin{document}
\begin{letter}{Guest Editors\\
  Special Issue ``Process Mining: Theory and Applications''\\
  \textit{Applied Sciences}}

\opening{Dear Guest Editors,}

We submit our manuscript, ``Reusable Neural Guidance for Petri Net Alignments
with Replay Verification'', by Alessandro Berti and Wil M. P. van der Aalst,
for consideration as a research article in this special issue of
\textit{Applied Sciences}.

Alignments provide executable explanations of deviations between event traces
and process models, but exact computation can become expensive. Our method,
LARA (Learned Alignment with Replay Assurance), uses one pretrained neural
scorer to guide alignment construction across Petri nets and activity
vocabularies. It combines learned move preferences with a decoder that tracks
the marking, independent replay,
and optional exact certification. Successful replay establishes executability
and a cost upper bound; completed exact search establishes optimality or
provides a replacement alignment. The paper also presents a synthetic
supervision procedure that pairs equivalent process representations with shared
observations and exact alignment targets.

On 512 synthetic test inputs, all candidates pass replay, and learned guidance
increases the proportion with optimal cost from 86.9\% to 91.0\% compared with
the same decoder without learned guidance. The largest improvement concerns
choices between transitions with duplicate activity labels. Experiments on
larger synthetic models and two public event logs examine reuse and the tradeoff
between candidate latency and quality. The evaluation also documents the
method's limits: established approximations generally produce better candidates,
and the current certification mode invokes exact search for every trace.

The manuscript fits the special issue by connecting process mining theory,
machine learning, and practical conformance analysis with explicit guarantees
for alignment results. The public implementation, synthetic
corpus, trained checkpoint, and experiment records support reproducibility;
an interactive workbench is available at \url{https://lara-pm.app}.
The manuscript discloses Wil M. P. van der Aalst's Celonis affiliation and the
use of generative AI tools.

% Author declarations: retain after confirming that both statements are accurate.
We confirm that neither this manuscript nor any part of its content has been
published in another journal or is being considered for publication elsewhere.
Both authors have approved the manuscript and agree to its submission to
\textit{Applied Sciences}.

Thank you for considering our manuscript.

\closing{Sincerely,}

\end{letter}
\end{document}
